\section{The question and the answer}
\label{sec:introduction}

A locally profinite space has a basis of compact-open sets. Locally, such
sets are particularly well behaved for sheaf cohomology: sections on a
compact open have no higher derived functors. Nevertheless, assembling
many compact opens can produce nonzero cohomology, and even nonzero
products of positive-degree classes. How large must the space be for
this to happen?

In \cite[Theorem C]{Aoki}, Aoki constructs, for each \(n\geq0\), a locally
profinite space \(U\) of cardinality \(\aleph_{2n+1}\) with classes
\(\eta_0,\ldots,\eta_n\in H^1(U;\F)\) such that

\[
\eta_0\smile\cdots\smile\eta_n\neq0.
\]

He then asks the following question.

\begin{quote}
\textbf{Aoki's Question 2.14.}
Is Theorem C optimal in terms of the weight (or cardinality) of \(U\)?
\end{quote}

We prove that the sharp bound for a product of \(n+1\) degree-one classes
is \(\aleph_{n+1}\), for both cardinality and weight. All factors in the
product can in fact be the same class. To keep the cohomological degree
visible, we write \(r=n+1\) throughout the proof.

A \emph{locally profinite Hausdorff space} will mean a locally compact,
Hausdorff, totally disconnected space. A compact Hausdorff totally
disconnected space is called profinite. Our examples are a profinite
space with one point removed, so they also satisfy the convention used
in \cite{Aoki}. The weight \(w(V)\) is the least cardinality of a basis for the topology.
All cohomology is ordinary sheaf cohomology; a constant
coefficient such as \(\F\) means the associated constant sheaf. Powers
of a degree-one class are cup powers. We work with the axiom of choice;
no continuum hypothesis is assumed.

\begin{theorem}\label{thm:main}
For every integer \(r\geq1\), there exist a locally profinite Hausdorff
space \(U_r\) and a class \(\eta_r\in H^1(U_r;\F)\) such that

\[
|U_r|=w(U_r)=\aleph_r,
\qquad \eta_r^r\neq0,
\qquad \eta_r^{r+1}=0.
\]

Moreover, for every abelian sheaf \(\mathcal M\) on \(U_r\),

\[
H^q(U_r;\mathcal M)=0\qquad(q>r).
\]

Conversely, for every locally profinite Hausdorff space \(V\) and every
abelian sheaf \(\mathcal M\) on \(V\),

\[
H^r(V;\mathcal M)\neq0
\quad\Longrightarrow\quad
|V|\geq\aleph_r\quad\text{and}\quad w(V)\geq\aleph_r.
\]

Thus \(\aleph_r\) is both the least cardinality and the least weight of a
locally profinite Hausdorff space supporting a nonzero product of \(r\)
degree-one constant-\(\F\) classes.\leanmark{main}
\end{theorem}

The theorem shows that Aoki's bound is optimal when \(n=0\), and can be
lowered for every \(n\geq1\). It also gives examples of cohomological
dimension exactly \(r\), where dimension is taken over all abelian sheaf
coefficients.

\subsection{Reading the paper and the formalization}

The proof is accompanied by a Lean development. Small references such as
\leanmark{main} point to the source concordance in
\cref{app:lean}, where each mathematical step is matched with explicit
files, declarations, and line numbers. These references can be skipped
on a first reading. The mathematical argument does not require knowledge
of Lean. \Cref{sec:formalization} explains precisely what is checked and
how to reproduce the verification.

\subsection{The proof in five steps}
\label{subsec:strategy}

There are two logically independent tasks: attaining the size
\(\aleph_r\), and excluding every smaller size. The first task contains
the main construction. Its guiding idea is to build a cocycle for which
being a boundary would force an impossible equation between functions.

\begin{enumerate}
\item \textbf{Put symmetry on coordinates of increasing sizes.}
Take \(r+1\) coordinates with underlying sizes
\(\aleph_0,\ldots,\aleph_r\). Each coordinate has a bit in \(\F\)
and a point at infinity. Identify points that differ by flipping every
bit, and remove the all-infinite point. This is \(U_r\). It has a cover
by \(r+1\) opens, according to which coordinate is finite.

\item \textbf{Write the class and its power explicitly.}
On the \(i\)-th open, choose the representative whose \(i\)-th bit is
zero. On an overlap, the difference of the two choices is
\(z_{ij}=s_i+s_j\). These differences give \(\eta_r\), and its top
power is represented by

\[
F_r=\prod_{j=0}^{r-1}(s_j+s_{j+1}).
\]

The intersections in the cover are disjoint unions of compact opens.
This makes their free sheaves projective, so an explicit finite
resolution computes ordinary cohomology and its cup product.

\item \textbf{Extract the restrictions on a possible boundary.}
If \(F_r\) were a boundary, it would be a sum
\(f_0+\cdots+f_r\). Each \(f_i\) would be invariant under the
simultaneous bit flip. It would also be constant outside finitely many
points in its \(i\)-th coordinate, with the other coordinates fixed:
this is exactly what continuity at infinity forces.

\item \textbf{Use the last coordinate to lower the degree.}
There are only \(\aleph_{r-1}\) configurations of the earlier
coordinates. We can therefore choose one base point in the last
coordinate that avoids every finite exceptional set for \(f_r\).
Adding the values at its two bits kills \(f_r\) and changes \(F_r\)
into \(F_{r-1}\). The other summands retain their symmetry and
almost-constancy. Induction reaches a two-coordinate contradiction:
an almost-constant function cannot differ by \(1\) from its bit flip
everywhere.

\item \textbf{Bound the cohomology of every smaller space.}
Independently, a cover by at most \(\aleph_n\) compact opens forces
cohomology to vanish above degree \(n\). Countable covers can be
disjointified; passing from one aleph to the next uses a continuous
transfinite filtration and raises the projective-dimension bound by
at most one. Both cardinality and weight control the size of a
compact-open cover. This proves the lower bounds.
\end{enumerate}

The roles of the three ingredients in the nonvanishing argument are
worth keeping distinct. \emph{Continuity} produces finite exceptional
sets. \emph{Cardinal growth} permits a choice avoiding all of them at
once. \emph{Symmetry} makes the two evaluations fit together into an
induction. The projective resolution connects this elementary argument
to sheaf cohomology.

\subsection{The first example as a guide}

For \(r=1\), the two coordinate sets have sizes \(\aleph_0\) and
\(\aleph_1\), and the cocycle is simply \(s_0+s_1\). A hypothetical
boundary would split it into two invariant functions, the first
almost constant in coordinate zero and the second almost constant in
coordinate one. Countably many finite exceptional sets in the second
coordinate cannot exhaust its uncountable set of base points. Evaluate
at both bits of a base point outside their union. The second function
cancels, and the resulting slice of the first would have to change by \(1\)
under bit flip.
Far enough from its own finite exceptional set, however, it is constant.
This contradiction is the base case for every higher cup power.

We now carry out the construction in \cref{sec:geometry}, the
cohomological calculation in \cref{sec:resolution,sec:cup}, and the
nonvanishing argument in \cref{sec:nonvanishing}. The independent
optimality proof occupies \cref{sec:optimality}.
